%% AUTO-GENERATED by scripts/make_kernel_results_tex.py -- do not edit by hand.
%% Every number here is re-derived from data/marlowe_h100 at build time.
\begin{table}[h]\centering\small
\caption{T1-kernel. Solve-rate alone is a compliance metric: a submission that ignores the
kernel instruction and rewrites the reference in torch verifies trivially. Reported instead as
a cascade. $^{\dagger}$ marks an incomplete cell. Generation budgets differ: the 0.5B--14B cells ran at $k{=}12$ and 32B at $k{=}6$, so 32B contributes half the candidates. The reported quantity is a \emph{rate} and is budget-independent in expectation, but 32B's estimate is correspondingly noisier. The curve is a \textbf{U}: it falls from 0.5B to a minimum of $0$ at 14B and recovers to $3.0\%$ at 32B (Fisher $p{=}0.015$ against 14B). Extraction policy: \texttt{strict}. Under \texttt{strict} the extractor discards 66\% of 0.5B generations and 5\% of 14B generations, so any trend across scale must be shown under both policies before it is attributed to the models.}
\begin{tabular}{@{}lrrrrr@{}}
\toprule
Scale & Generations & Parsed & Attempted kernel & Verified & Kernel-verified \\
\midrule
0.5B & 348 & 118 (34\%) & 74 (63\%) & 11 & 4 (5.4\%) \\
1.5B & 348 & 132 (38\%) & 74 (56\%) & 11 & 3 (4.1\%) \\
3B & 348 & 253 (73\%) & 184 (73\%) & 9 & 6 (3.3\%) \\
7B & 348 & 335 (96\%) & 262 (78\%) & 5 & 4 (1.5\%) \\
14B & 348 & 331 (95\%) & 219 (66\%) & 0 & 0 (0.0\%) \\
32B & 174 & 173 (99\%) & 167 (97\%) & 5 & 5 (3.0\%) \\
\bottomrule
\end{tabular}
\end{table}

\begin{table}[h]\centering\small
\caption{T1-kernel robustness to extraction policy. Verify-given-attempt under
\texttt{KA\_EXTRACT=strict} and \texttt{lenient}; identical tasks, $k$, and token budget,
differing only in how a generation is converted to a candidate. The endpoint contrast
(0.5B vs 14B) holds under both policies and strengthens under \texttt{lenient}. Note that 14B is the curve's \emph{minimum} and not its largest scale: the completed 32B sweep recovers to $3.0\%$ (Fisher $p{=}0.015$ against 14B), so this is a contrast with the floor. The monotone ordering does not: it is an artifact of a filter whose severity correlates with scale.}
\begin{tabular}{@{}lrrrrrr@{}}
\toprule
& \multicolumn{3}{c}{\texttt{strict}} & \multicolumn{3}{c}{\texttt{lenient}} \\
\cmidrule(lr){2-4}\cmidrule(lr){5-7}
Scale & Attempts & Verified & Rate & Attempts & Verified & Rate \\
\midrule
0.5B & 74 & 4 & 5.41\% & 254 & 24 & 9.45\% \\
1.5B & 74 & 3 & 4.05\% & 237 & 7 & 2.95\% \\
3B & 184 & 6 & 3.26\% & 328 & 26 & 7.93\% \\
7B & 262 & 4 & 1.53\% & 344 & 10 & 2.91\% \\
14B & 219 & 0 & 0.00\% & 338 & 2 & 0.59\% \\
\bottomrule
\end{tabular}
\end{table}

\begin{table}[h]\centering\small
\caption{14B verify-given-attempt, replicated across independent routes. Two runs differing in route, sampling budget $k$, and teacher bank. The model parses as a kernel attempt in 98\% of its candidates and almost never produces one that verifies, so the result is not an artifact of extraction, prompt compliance, or a single unlucky run.}
\begin{tabular}{@{}llrrr@{}}
\toprule
Run & Route & $k$ & Attempted & Kernel-verified \\
\midrule
\texttt{t1kl\_q14} & T1 (lenient) & 12 & 338 & 2 (0.59\%) \\
\texttt{r4\_kernel\_q14} & R4 & 6 & 171 & 1 (0.58\%) \\
\bottomrule
\end{tabular}
\end{table}

\begin{table}[h]\centering\small
\label{tab:nonllm}
\caption{A non-LLM baseline on the benchmark's own terms: \texttt{torch.compile(mode=\"max-autotune\")} scored exactly as a model submission is (speedup over the compiled baseline, same per-task roofline ceiling, same timing harness). It lands \emph{below} correct-at-parity because it is slower than default \texttt{torch.compile} on 27 of 29 tasks. A model scoring 0.50 therefore beats production autotuning here. Reported per tier because compile gain over eager differs by ${\sim}12\times$ between L1 and L2.}
\begin{tabular}{@{}lrrr@{}}
\toprule
Tier & Tasks & Median score & Median speedup vs compiled \\
\midrule
L1 & 10 & 0.381 & 0.611 \\
L2 & 15 & 0.408 & 0.868 \\
L3 & 4 & 0.505 & 1.063 \\
\midrule
All & 29 & \textbf{0.408} & 0.865 \\
\bottomrule
\end{tabular}
\end{table}

\begin{table}[h]\centering\small
\label{tab:passrate}
\caption{T2 weight-RSI under pass rate ($n_{\mathrm{ok}}/k$), the \texttt{h100/passrate} set. Reported trajectory-level: one value per cell, the mean of its last two rounds, since rounds inside a cell share an adapter. Lineage beats a matched reset learner that sees identical per-round data, which isolates weight accumulation. Teacher injection changes it by +0.005. \textbf{No interval is quoted} at $n{=}2$ trajectories per arm. $C_{\mathrm{lineage}}$ reaches the metric's ceiling of 1.0 in three of four cells, so the final rounds measure the bound rather than the learner. Per PREREGISTRATION Amendment 1 this board is never pooled with, or compared against, any headroom result.}
\begin{tabular}{@{}llrrr@{}}
\toprule
Cell & Arm & Rounds & lineage $-$ reset & final $C_{\mathrm{lineage}}$ \\
\midrule
\texttt{control\_q15\_s1} & control & 8 & +0.450 & 0.88 \\
\texttt{control\_q15\_s2} & control & 8 & +0.430 & 0.98 \\
\texttt{inject\_q15\_s1} & inject & 8 & +0.500 & 1.00 \\
\texttt{inject\_q15\_s2} & inject & 8 & +0.390 & 1.00 \\
\midrule
\multicolumn{3}{@{}l}{control, trajectory mean} & \textbf{+0.440} & \\
\multicolumn{3}{@{}l}{inject, trajectory mean} & \textbf{+0.445} & \\
\bottomrule
\end{tabular}
\end{table}

\begin{table}[h]\centering\small
\label{tab:contrast}
\caption{The same experiment under two scorers. \texttt{t2kc} and \texttt{t2kp} differ in \texttt{KA\_SCORE} alone: same model, seed, $k$, lab, task bank and arms. $\dagger$ marks a round where \emph{both} arms have reached correct-at-parity ($\geq 0.49$). In all 10 such rounds the headroom contrast is within $0.004$ of zero (mean $-0.0002$), while the same rounds report a mean of $+0.436$ under pass rate. The contrast measures the \emph{gap} between the arms and nothing else: across all four headroom cells and both scales, the 18 rounds with arms within $0.02$ of each other average $-0.0003$ and never exceed $0.004$, while the 14 rounds with a gap average $+0.144$. It is not broken by saturation and recovers when a gap reopens (\texttt{q3\_s2} reads $+0.002$ at $0.50/0.50$ then $+0.201$ at $0.50/0.30$). The difficulty is that best-of-$k$ compresses both arms onto correct-at-parity and holds them there, and a matched reset learner arrives within three or four rounds, so the measurement's useful lifetime is fixed by the control arm rather than by the treatment. The two boards are never \emph{pooled} (Amendment 1); they are placed side by side here to compare the scorers, not the results.}
\begin{tabular}{@{}lrrrr@{}}
\toprule
Cell & Round & $C_{\mathrm{lin}}$ / $C_{\mathrm{reset}}$ & headroom & pass rate \\
\midrule
q15\_s1 & 0 & 0.20 / 0.10 & +0.101 & +0.020 \\
q15\_s1 & 1 & 0.40 / 0.20 & +0.200 & +0.080 \\
q15\_s1 & 2 & 0.51 / 0.20 & +0.304 & +0.220 \\
q15\_s1 & 3 & 0.50 / 0.51 & -0.004$\dagger$ & +0.440 \\
q15\_s1 & 4 & 0.50 / 0.50 & +0.000$\dagger$ & +0.520 \\
q15\_s1 & 5 & 0.50 / 0.50 & -0.000$\dagger$ & +0.380 \\
q15\_s1 & 6 & 0.50 / 0.50 & +0.001$\dagger$ & +0.380 \\
q15\_s1 & 7 & 0.50 / 0.50 & -0.001$\dagger$ & +0.520 \\
q15\_s2 & 0 & 0.00 / 0.20 & -0.201 & +0.020 \\
q15\_s2 & 1 & 0.29 / 0.10 & +0.189 & -0.020 \\
q15\_s2 & 2 & 0.30 / 0.30 & -0.004 & +0.080 \\
q15\_s2 & 3 & 0.50 / 0.50 & +0.001$\dagger$ & +0.380 \\
q15\_s2 & 4 & 0.50 / 0.50 & -0.000$\dagger$ & +0.360 \\
q15\_s2 & 5 & 0.50 / 0.50 & +0.001$\dagger$ & +0.520 \\
q15\_s2 & 6 & 0.50 / 0.50 & -0.000$\dagger$ & +0.400 \\
q15\_s2 & 7 & 0.50 / 0.50 & +0.000$\dagger$ & +0.460 \\
\bottomrule
\end{tabular}
\end{table}

\begin{table}[h]\centering\small
\label{tab:fed}
\caption{The same lab and the same contrast, with the learner fed. These cells run \texttt{--n-train 35} against a verified kernel bank; the registered primary runs \texttt{--n-train 3}, which the throughput arithmetic puts at $0.96$ examples per round. Mean $n_{\mathrm{ex}}$ here is 50.8 and 27.4, so both trajectories clear Amendment 6's floor of two by more than a factor of ten, and both carry \texttt{adapter\_restored}. Every round of both is reported. \textbf{No interval is quoted and no effect is claimed}: $n=2$ trajectories against a registered minimum of three seeds, and a 95\% interval on a two-point mean would assert precision the design cannot support. The trajectory-level mean is $+0.101$. The column that carries weight is $n_{\mathrm{ex}}$, which grows within both runs (6 to 104 and 7 to 84) and is the input the contrast lacked everywhere else in this report. Registered rule: the 95\% interval excludes zero ($+0.101\,[+0.018,+0.183]$, $n=5$) and mean(last-2) exceeds $0.05$ in every cell. The rule requires \emph{both}, so the result \textbf{holds}.}
\begin{tabular}{@{}lrrlrr@{}}
\toprule
Cell & Rounds & mean $n_{\mathrm{ex}}$ & \texttt{self}$-$\texttt{fresh} per round & mean & last 2 \\
\midrule
\texttt{q15\_s1} & 5 & 50.8 & $+0.074$ $+0.132$ $+0.206$ $+0.130$ $+0.089$ & $+0.126$ & $+0.110$ \\
\texttt{q15\_s2} & 5 & 27.4 & $-0.001$ $-0.088$ $+0.073$ $+0.135$ $+0.131$ & $+0.050$ & $+0.133$ \\
\texttt{q15\_s3} & 5 & 55.0 & $+0.033$ $+0.125$ $+0.272$ $+0.193$ $+0.161$ & $+0.157$ & $+0.177$ \\
\texttt{q15\_s4} & 5 & 20.8 & $-0.021$ $-0.023$ $-0.063$ $+0.072$ $+0.094$ & $+0.012$ & $+0.083$ \\
\texttt{q15\_s5} & 5 & 59.0 & $+0.117$ $+0.213$ $+0.243$ $+0.148$ $+0.075$ & $+0.159$ & $+0.111$ \\
\bottomrule
\end{tabular}
\end{table}

\begin{table}[h]\centering\small
\label{tab:dose}
\caption{Throughput as a dose. Every cell here runs the same bank, model, $k$, arms and code path, and differs only in \texttt{--n-train}, so the verified examples the learner receives is an independent variable rather than a difference between two configurations. The \texttt{fed} cells reported elsewhere changed the bank \emph{and} \texttt{n\_train} together and cannot separate the two. $\dagger$ marks a rung whose mean $n_{\mathrm{ex}}$ falls below the floor of two set by PREREGISTRATION Amendment 6, at which a trajectory is reported as starved rather than as a measurement of the contrast. That the low rungs are starved is the \emph{result} here and not grounds for excluding them: it is what varying the dose was meant to show. A starved rung's contrast is set by whichever arm happened to receive data, so its sign carries no information about self-training.}
\begin{tabular}{@{}lrrlr@{}}
\toprule
\texttt{n\_train} & Cells & mean $n_{\mathrm{ex}}$/round & per-trajectory mean & pooled \\
\midrule
3 & 2 & 0.20$^{\dagger}$ & $+0.012$, $-0.017$ & $-0.003$ \\
10 & 2 & 0.60$^{\dagger}$ & $-0.124$, $+0.002$ & $-0.061$ \\
35 & 2 & 39.10 & $+0.126$, $+0.050$ & $+0.088$ \\
\bottomrule
\end{tabular}
\end{table}

\begin{table}[h]\centering\small
\label{tab:kl}
\caption{Regularized RSI. Each row trains under $\theta(t{+}1) = \arg\max \mathbb{E}[R] - \lambda\,\mathrm{KL}(\pi(\theta)\,\|\,\pi(\theta_0))$, with the reference policy obtained by disabling the LoRA adapter rather than loading a second model. $\lambda = 0$ is the unregularized objective and is the same code path the \texttt{fed} cells ran. The realized KL is reported beside the setting because a $\lambda$ that does not move the divergence is not a treatment. lambda=2 (0 of 2), lambda=5 (1 of 2), lambda=20 (1 of 2) not yet at depth and omitted.}
\begin{tabular}{@{}lrrrrr@{}}
\toprule
$\lambda$ & Cells & realized KL & mean $n_{\mathrm{ex}}$ & \texttt{self}$-$\texttt{fresh} & last 2 \\
\midrule
$0$ & 2 & -- & 39.1 & $+0.088$ & $+0.121$ \\
$0.05$ & 2 & 0.0047 & 42.1 & $+0.118$ & $+0.180$ \\
$0.2$ & 2 & 0.0033 & 40.8 & $+0.111$ & $+0.116$ \\
$1.0$ & 2 & 0.0023 & 23.6 & $+0.105$ & $+0.121$ \\
\bottomrule
\end{tabular}
\end{table}

\begin{table}[h]\centering\small
\label{tab:saturation}
\caption{A confident null, manufactured by a one-token change. \texttt{nb20} and \texttt{bk20} differ only in \texttt{KA\_KERNEL\_BANK}; model, seed, rounds, $k$, $n_{\mathrm{train}}=20$ and $n_{\mathrm{held}}=9$ are identical. Rounds are split on whether the frozen control has already saturated the scorer ($C_{\mathrm{reset}} \ge 0.49$; a correct kernel at parity scores exactly $0.50$), not on which bank they came from. Saturated rounds measure a mean absolute contrast of $0.0005$ against $0.1253$ where the instrument has room, a factor of $251$. The split holds \emph{within} \texttt{nb20}, and \texttt{nb20}'s unsaturated rounds agree with \texttt{bk20}'s, so the task bank does not change the effect --- it changes how often the instrument can resolve it. On its own, \texttt{nb20} reads as a tight equivalence result.}
\begin{tabular}{@{}llrrr@{}}
\toprule
Cell & Bank & saturated & mean $|\Delta|$ saturated & mean $|\Delta|$ has room \\
\midrule
\texttt{nb20} & default, 29 tasks & 8/16 & $0.0005$ & $0.1129$ \\
\texttt{bk20} & kernel, 124 tasks & 0/16 & --- & $0.1316$ \\
\midrule
pooled & --- & 8/32 & $0.0005$ & $0.1253$ \\
\bottomrule
\end{tabular}
\end{table}

\begin{table}[h]\centering\small
\label{tab:depth}
\caption{The same configuration read at two depths. \texttt{deep\_q15} and \texttt{fed\_q15} differ in one token of the command line, \texttt{--rounds 15} against \texttt{--rounds 5}; model, seed, $k$, $n_{\mathrm{train}}$, task bank and GPU map are identical. Every seed is positive over rounds 1--5 and negative over rounds 14--15, so stopping at five rounds and stopping at fifteen support opposite conclusions about whether self-training beats a frozen control. The loop is not starved at any point --- $n_{\mathrm{ex}}$ stays far above the starvation floor of two --- and round-0 retention falls to as little as $0.27$, which is the mechanism we attribute the reversal to (Section~\ref{sec:instrument}).}
\begin{tabular}{@{}lrrrrr@{}}
\toprule
Cell & Rounds & mean $\Delta$ r1--5 & mean $\Delta$ r14--15 & final retention & mean $n_{\mathrm{ex}}$ \\
\midrule
\texttt{q15\_s1} & 15 & $+0.180$ & $-0.183$ & 0.27 & 85 \\
\texttt{q15\_s2} & 15 & $+0.084$ & $-0.270$ & 0.26 & 69 \\
\texttt{q15\_s3} & 15 & $+0.082$ & $-0.089$ & 0.43 & 75 \\
\bottomrule
\end{tabular}
\end{table}

\begin{table}[h]\centering\small
\label{tab:lifetime}
\caption{When each arm reaches correct-at-parity, and what happens to the contrast afterwards. The \emph{gap} between the two arrival rounds bounds the largest contrast the design can ever observe, which is a property of the design rather than of the model. Where the arms arrive together the contrast never exceeds $0.062$ at any point in the run, however much data the learner receives: the \texttt{fedoc} cells take 350 verified examples per round and peak below that. Where they arrive two rounds apart it reaches $0.15$ to $0.27$. The \emph{final} column is near zero in every cell whose control has arrived, whatever the gap, because by the last round both arms sit at the ceiling -- so the gap caps how large an effect is visible, not merely when it stops being visible. Reported per cell and never pooled: the arrival rounds are integers on different models, and their mean would describe a cell that does not exist.}
\begin{tabular}{@{}lrrrrrrr@{}}
\toprule
Cell & Rounds & treat.\ at & ctrl.\ at & gap & peak & final & max $n_{\mathrm{ex}}$ \\
\midrule
\texttt{fed\_q05\_s1} & 5 & 2 & 4 & 2 & $-0.152$ & $+0.005$ & 275 \\
\texttt{fed\_q05\_s2} & 5 & 3 & -- & $>1$ & $+0.140$ & $+0.004$ & 257 \\
\texttt{fed\_q05\_s3} & 5 & 1 & -- & $>3$ & $+0.017$ & $+0.005$ & 304 \\
\texttt{fed\_q15\_s1} & 5 & 4 & -- & $>0$ & $+0.206$ & $+0.089$ & 104 \\
\texttt{fed\_q3\_s1} & 5 & 2 & 2 & 0 & $+0.042$ & $+0.014$ & 241 \\
\texttt{fed\_q3\_s2} & 5 & 2 & 4 & 2 & $+0.017$ & $+0.017$ & 194 \\
\texttt{fed\_q3\_s3} & 5 & 2 & 2 & 0 & $+0.029$ & $+0.006$ & 219 \\
\texttt{fedds\_s1} & 5 & 2 & 4 & 2 & $-0.272$ & $+0.002$ & 346 \\
\texttt{fedds\_s2} & 5 & 2 & 4 & 2 & $+0.183$ & $+0.011$ & 307 \\
\texttt{fedds\_s3} & 5 & 2 & 4 & 2 & $-0.136$ & $-0.002$ & 327 \\
\texttt{fedoc\_s1} & 5 & 1 & 1 & 0 & $+0.062$ & $-0.002$ & 350 \\
\texttt{fedoc\_s2} & 5 & 1 & 1 & 0 & $-0.006$ & $-0.002$ & 348 \\
\texttt{fedoc\_s3} & 5 & 1 & 1 & 0 & $-0.039$ & $-0.004$ & 348 \\
\bottomrule
\end{tabular}
\end{table}

\begin{table}[h]\centering\small
\label{tab:collapse}
\caption{Verified self-training past round five. The treatment arm peaks and then loses capability \emph{absolutely}: the drop column is its final score minus its peak. Two columns rule out the alternatives within each run. \texttt{trainC} is the model's score on its own training split and falls alongside the held-out score, which overfitting would not do. $C_{\mathrm{ctrl}}$ is the round0-replay arm --- identical LoRA, learning rate, schedule and round count, trained on verified self-data \emph{frozen at round 0} --- and it does not collapse, so the cause is the recursion rather than repeated training. It also reaches a higher score than the treatment ever does, so nothing is capping the treatment. Every training example in both arms passed execution against an fp32 reference.}
\begin{tabular}{@{}lrrrrrrrrrr@{}}
\toprule
Cell & Rounds & peak at & peak $C$ & final $C$ & drop & $C_{\mathrm{ctrl}}$ max & final & \texttt{trainC} & contrast @5 & contrast final \\
\midrule
\texttt{deep\_q15\_s1} & 15 & 4 & 0.47 & 0.18 & $-0.29$ & 0.48 & 0.48 & 0.50$\to$0.27 & $+0.182$ & $-0.251$ \\
\texttt{deep\_q15\_s2} & 15 & 4 & 0.49 & 0.16 & $-0.33$ & 0.49 & 0.49 & 0.49$\to$0.17 & $+0.071$ & $-0.321$ \\
\texttt{deep\_q15\_s3} & 15 & 2 & 0.48 & 0.24 & $-0.24$ & 0.51 & 0.49 & 0.49$\to$0.25 & $+0.082$ & $-0.097$ \\
\texttt{deepp\_q15\_s3} & 6 & 4 & 0.42 & 0.32 & $-0.09$ & 0.23 & 0.23 & 0.47$\to$0.47 & $+0.245$ & $+0.185$ \\
\bottomrule
\end{tabular}
\end{table}
